\subsection{Thiết kế kho dữ liệu}
\subsubsection*{Bảng chiều (dimension)}
Kho dữ liệu có 10 bảng chiều. Mỗi bảng có khoá thay thế dạng số nguyên tự tăng, khoá tự nhiên ràng buộc duy nhất, và một dòng Unknown với khoá -1 để bảng fact không bao giờ có khoá rỗng. Chiều thời gian \code{dim\_\allowbreak{}date} phủ 1.461 ngày từ 2016 đến 2019 với ngày lễ Brazil và Black Friday, và được dùng lại cho nhiều mốc thời gian của đơn hàng qua 11 view đóng vai. Chiều địa lý \code{dim\_\allowbreak{}geography} có một dòng cho mỗi mã bưu chính, hợp nhất từ bảng toạ độ, khách hàng và người bán, với phân cấp Zip $\rightarrow$ Thành phố $\rightarrow$ Bang $\rightarrow$ Vùng. Chiều khách hàng \code{dim\_\allowbreak{}customer} áp dụng SCD loại 2: khi một khách đổi mã bưu chính qua các đơn, hệ thống đóng phiên bản cũ và mở phiên bản mới; trên dữ liệu đầy đủ có 96.096 khách và thêm 242 phiên bản (khoảng 0,25\%), nên chi phí lưu lịch sử thấp. Chiều người bán \code{dim\_\allowbreak{}seller} là chiều dùng chung (conformed) nối quy trình bán hàng với quy trình thu hút seller, dùng SCD loại 1; 462 seller chỉ có trong CRM vẫn có dòng với địa chỉ chưa biết. Chiều sản phẩm gom 73 danh mục vào 13 nhóm ngành hàng và phân loại cỡ hàng theo cân nặng. Các thuộc tính nhỏ của deal được gom vào một chiều rác (junk dimension) gồm 503 tổ hợp, và nhân viên sales được tách thành hai vai trò SDR và SR trong cùng một bảng. Toàn bộ 39 phép kiểm thử dimension đều đạt: mỗi chiều có đúng một dòng Unknown, mỗi khách có đúng một phiên bản hiện hành và các khoảng hiệu lực không chồng lấn.

\begin{table}[htbp]
\centering\small
\caption{Các bảng chiều: grain, số dòng (gồm dòng Unknown $-1$) và loại SCD}
\label{tab:dims}
\begin{tabular}{llrc}
\toprule
\textbf{Chiều} & \textbf{Grain} & \textbf{Số dòng} & \textbf{SCD} \\
\midrule
dim\_date & 1 ngày (2016-01-01 → 2020-12-31) & 1.828 & -- \\
dim\_geography & 1 mã bưu chính (zip) & 19.178 & 1 \\
dim\_customer & 1 phiên bản địa chỉ của 1 khách thật & 96.339 & 2 \\
dim\_seller & 1 seller (chiều dùng chung) & 3.558 & 1 \\
dim\_product & 1 sản phẩm & 32.952 & 1 \\
dim\_order\_status & 1 trạng thái đơn & 9 & 1 \\
dim\_payment\_type & 1 hình thức thanh toán & 6 & 1 \\
dim\_lead\_origin & 1 giá trị origin của MQL & 11 & 1 \\
dim\_lead\_profile & 1 tổ hợp thuộc tính deal (junk dimension) & 503 & -- \\
dim\_sales\_rep & 1 nhân viên và vai trò (SDR/SR) & 55 & -- \\
\bottomrule
\end{tabular}
\end{table}

\subsubsection*{Bảng fact và mô hình vật lý}
Kho dữ liệu có sáu bảng fact thuộc cả ba loại của Kimball. Hai bảng giao dịch ở mức chi tiết nhất là \code{fact\_\allowbreak{}order\_\allowbreak{}item} (một dòng cho mỗi sản phẩm trong đơn, 112.650 dòng khi nạp đủ) và \code{fact\_\allowbreak{}payment} (một dòng cho mỗi lần thanh toán); \code{fact\_\allowbreak{}review} ghi điểm đánh giá, kèm cờ cho biết đơn có giao trễ hay không. Bảng \code{fact\_\allowbreak{}order\_\allowbreak{}fulfillment} là snapshot tích luỹ: mỗi đơn một dòng và được cập nhật khi đơn đi qua từng mốc (mua, duyệt, bàn giao vận chuyển, giao khách, ngày dự kiến), từ đó tính thời gian từng chặng, số ngày trễ và số ngày Olist hứa giao. Quy trình thu hút seller dùng snapshot tích luỹ thứ hai \code{fact\_\allowbreak{}lead\_\allowbreak{}funnel} (mỗi khách tiềm năng một dòng, đi từ liên hệ đầu, chốt deal đến đơn đầu tiên), còn \code{fact\_\allowbreak{}seller\_\allowbreak{}monthly} là snapshot định kỳ, mỗi seller mỗi tháng một dòng, lưu dày đặc để đo được cả những tháng seller không bán. Mọi khoá ngoại đều không rỗng; nếu không tra được chiều thì dùng dòng Unknown và ghi vào nhật ký ETL. Khoá khách hàng luôn trỏ tới phiên bản địa chỉ có hiệu lực vào ngày mua. Các khoảng thời gian âm do lỗi nguồn được đặt rỗng và đếm lại, không xoá đơn. Sau khi nạp, tổng giá bán ở fact khớp từng đồng với dữ liệu nguồn (13.591.643,70 BRL khi nạp đủ), và 112 phép kiểm thử về số dòng, tổng tiền, grain, khoá ngoại và logic nghiệp vụ đều đạt. Ở mức vật lý, mỗi bảng fact có khoá chính theo grain, khoá ngoại sang các chiều và chỉ mục B-tree trên mọi khoá ngoại.

\begin{table}[htbp]
\centering\small
\caption{Các bảng fact: loại, grain và số dòng}
\label{tab:facts}
\begin{tabular}{lllr}
\toprule
\textbf{Fact} & \textbf{Loại} & \textbf{Grain} & \textbf{Số dòng} \\
\midrule
fact\_order\_item & Transaction & 1 dòng hàng của 1 đơn & 112.650 \\
fact\_order\_fulfillment & Accumulating snapshot & 1 đơn (5 mốc thời gian) & 99.441 \\
fact\_payment & Transaction & 1 lần thanh toán & 103.886 \\
fact\_review & Transaction & 1 review của 1 đơn & 99.224 \\
fact\_lead\_funnel & Accumulating snapshot & 1 MQL & 8.000 \\
fact\_seller\_monthly & Periodic snapshot (dày đặc) & 1 seller × 1 tháng & 36.583 \\
\bottomrule
\end{tabular}
\end{table}

\hinh{3-3-4_erd_dw.png}{Mô hình dữ liệu của kho (data model; bản nháp tự động)}{fig:erd-dw}{1.0}
\hinh{3-3-4_constellation.png}{Mô hình constellation theo data mart: fact (cam), dimension (xanh) (bản nháp tự động)}{fig:constellation}{0.9}
\todo{Yêu cầu của môn: \textbf{data model logic và vật lý} của staging / data mart / kho, dùng màu phân biệt fact (cam) và dimension (xanh), nhóm các vùng liên quan bằng khung. Vẽ lại bằng draw.io từ docs/diagrams/data-model-theo-mart.drawio. Thêm: hệ thống chiều khái niệm (phân cấp từng dimension), thống kê số bản ghi từng dimension (đã có Bảng \ref{tab:dims}), phân nhóm các dimension.}
\todo{Thêm bus matrix (docs/DW\_DESIGN.md mục 4) và sơ đồ phân cấp: ngày, địa lý, sản phẩm, seller, khách (SCD2).}
